\chapter{$\mathbf{R}^3$ 空间曲线坐标系中的向量分析}

\problem{2.1}{在柱坐标 $(\rho,\varphi,z)$ 下，试计算其坐标基 $\boldsymbol{e}_\rho,\boldsymbol{e}_\varphi,\boldsymbol{e}_z$（可称为空间中的标架场）的散度和旋度，即
\[
  \nabla\cdot\boldsymbol{e}_\rho,\quad \nabla\cdot\boldsymbol{e}_\varphi,\quad \nabla\cdot\boldsymbol{e}_z \quad\text{和}\quad \nabla\times\boldsymbol{e}_\rho,\quad \nabla\times\boldsymbol{e}_\varphi,\quad \nabla\times\boldsymbol{e}_z
\]
所得的计算结果，能否说明它们的物理意义？}
\answerarea

\problem{2.2}{在球坐标 $(r,\theta,\varphi)$ 下，试计算其坐标基矢 $\boldsymbol{e}_r,\boldsymbol{e}_\theta,\boldsymbol{e}_\varphi$ 所构成的标架场的散度和旋度，即
\[
  \nabla\cdot\boldsymbol{e}_r,\quad \nabla\cdot\boldsymbol{e}_\theta,\quad \nabla\cdot\boldsymbol{e}_\varphi \quad\text{和}\quad \nabla\times\boldsymbol{e}_r,\quad \nabla\times\boldsymbol{e}_\theta,\quad \nabla\times\boldsymbol{e}_\varphi
\]
所得的计算结果能否用物理学中的某些场量给予恰当的解释？}
\answerarea
\clearpage

\problem{2.3}{若有一个力场在直角坐标系下的表达式为
\[
  \boldsymbol{F} = -\frac{y}{x^2+y^2}\boldsymbol{e}_x + \frac{x}{x^2+y^2}\boldsymbol{e}_y
\]
\begin{enumerate}
  \item 请用柱坐标系把此力场表示出来；
  \item 在柱坐标系下计算此力场的散度和旋度；
  \item 若有一物体在此力场中沿 $\boldsymbol{e}_\rho$ 方向运动一单位长度时，求此力场 $\boldsymbol{F}$ 所做的功；
  \item 若有一物体在此力场的作用下，沿 $x$-$y$ 平面的一单位圆顺时针运动一圈，求此力场 $\boldsymbol{F}$ 所做的功，并分析此功与力场 $\boldsymbol{F}$ 的旋度的关系。
\end{enumerate}}
\answerarea

\problem{2.4}{请证明
\begin{enumerate}
  \item $\nabla r^n = n r^{n-2}\boldsymbol{r}$；
  \item $\nabla^2 r^n = n(n+1) r^{n-2}$，
\end{enumerate}
其中 $r=|\boldsymbol{r}|$。}
\answerarea
\clearpage

\problem{2.5}{假设点电荷 $q$ 的势为
\[
  \psi(\boldsymbol{r}) = \frac{q}{r^{1+\varepsilon}},\quad 0<\varepsilon<1
\]
在球坐标系下，计算此势（在 $r\neq 0$ 处）的场强 $\boldsymbol{E}\ (\boldsymbol{E}=-\nabla\psi)$ 并计算：
\begin{enumerate}
  \item $\nabla\cdot\boldsymbol{E}$；
  \item $\nabla\times\boldsymbol{E}$；
  \item 对半径为 $R$ 的球形区域，推导类似 Gauss 定律的结果，对于此结果，检验当 $\varepsilon\to 0$ 时的极限情况，并同点电荷的 Coulomb 势 $\psi(\boldsymbol{r})=\dfrac{q}{r}$ 做比较，说明这一电势的改变所引起的物理结果的变化。
\end{enumerate}}
\answerarea

\problem{2.6}{在球坐标系下，已知向量场 $\boldsymbol{A}(\boldsymbol{r})$ 的旋度为
\[
  \nabla\times\boldsymbol{A} = \frac{1}{r^2}\boldsymbol{e}_r
\]
求此向量场 $\boldsymbol{A}$。所得的场 $\boldsymbol{A}$ 是唯一的吗？如果不唯一，请说明理由。}
\answerarea
\clearpage
